\section{Shielded value on Flux}
\label{sec:shielded-value}\label{sec:privacy}

\fluxd is a Zcash fork, and its privacy layer is not something Flux built: it is Zcash's Sprout and
Sapling shielded-transaction machinery, carried forward through an intermediate ``Zel'' rebrand,
compiled into the node today, and only partially wired up for Flux's own chain. Every subsection
below states, component by component, whether what is described is inherited from Zcash unmodified,
inherited and then modified by Flux, or original to Flux. Almost nothing in this section is the last
category: the cryptography is Zcash's, and Flux's own work here is mostly a set of consensus-level
switches layered on top of it — some left on, several left off.

\fluxd's consensus rules and wallet support two address and transaction types side by side today:
transparent (\taddr) and shielded (\zaddr) (\shipped, \srccode{fluxd/\allowbreak{}src/\allowbreak{}primitives/\allowbreak{}transaction.h}{684-686,967-969}).
\textbf{That coexistence is ending.} Earlier drafts of this paper recorded it as permanent, on the
strength of three consistent statements: ``Optional at the protocol level, never mandatory. A
transparent layer stays permanently'' \srcroadmap{flux-roadmap-public.md:171}; the non-goals list
ruling out ``Making privacy mandatory'' for the same reason
\srcroadmap{flux-roadmap-public.md:206}; and the design intake's ``Both payment types stay,
permanently'' \srcintent{design intake}. Those statements are accurately reported and were
accurate when made. They have been superseded: the team has decided to \textbf{retire both shielded
pools} \srcintent{design intake} (\planned). The transparent layer stays; the
shielded layer does not.

This section is therefore written in two registers, and the reader should know which is which. Every
subsection up to and including \S\ref{sec:rpc-surface} describes the machinery \emph{as deployed
today} --- \shipped, present tense, accurate at the time of writing, and live on mainnet holding real
value. \S\ref{sec:shielded-retirement} states the decision to remove it, the measurement that
motivated it, and the schedule. The descriptive half is retained in full rather than cut: a reference
description of Flux that omitted a mechanism currently running would be incomplete, and a retirement
plan is only auditable against a precise account of what is being retired.

\subsection{Lineage: Bitcoin, Zcash, Zel, Flux}
\label{sec:shielded-lineage}

The in-repository license chain states the lineage plainly: \texttt{COPYING} carries copyright
headers running Bitcoin Core developers (2009--2019) $\to$ Zcash developers (2016--2019) $\to$ Zel
developers (2017--2019) $\to$ Flux Developers (2018--2022) (\shipped,
\srccode{fluxd/\allowbreak{}COPYING}{1-5}). That is the authoritative lineage as the code itself records it, and
it is stated here as a fact about provenance, not as an apology: a large fraction of \fluxd is
inherited Zcash code, and the paper treats every claim below accordingly.

The renaming from Zcash's \texttt{libzcash} to Flux's \texttt{libflux} was not applied uniformly.
Most files under \texttt{src/flux/} carry only a Flux copyright line (e.g.\
\srccode{fluxd/\allowbreak{}src/\allowbreak{}flux/\allowbreak{}Note.hpp}{1}, \srccode{fluxd/\allowbreak{}src/\allowbreak{}flux/\allowbreak{}NoteEncryption.hpp}{1},
\srccode{fluxd/\allowbreak{}src/\allowbreak{}flux/\allowbreak{}Proof.hpp}{1}, \srccode{fluxd/\allowbreak{}src/\allowbreak{}flux/\allowbreak{}Address.hpp}{1},
\srccode{fluxd/\allowbreak{}src/\allowbreak{}flux/\allowbreak{}JoinSplit.hpp}{1}), while \srccode{fluxd/\allowbreak{}src/\allowbreak{}flux/\allowbreak{}zip32.h}{1-2} retains both
the original Zcash and the added Flux copyright line — the header rewrite was partial. Under
\texttt{src/}, 395 files carry \texttt{"The Flux Developers"} and 24 still contain the literal string
\texttt{"Zelcash"}; repository-wide, 58 files still carry \texttt{"The Zcash developers"} with no Flux
line added (\shipped, \srccode{fluxd}{repo-wide grep}). The historical header \texttt{src/\allowbreak{}flux/\allowbreak{}Zelcash.\allowbreak{}h}, which
defines the shielded protocol's size constants (Table~\ref{tab:shielded-components}), still bears that
name and is directly \texttt{\#include}d by \texttt{Note.hpp}, \texttt{Note\allowbreak{}Encryption.\allowbreak{}hpp} and
\texttt{Address.hpp} (\shipped, \srccode{fluxd/\allowbreak{}src/\allowbreak{}flux/\allowbreak{}Note.hpp}{9}). Other rebrand artifacts persist
unmodified: the standalone consensus library's pkg-config template still emits
\texttt{-\allowbreak{}lzelcashconsensus} and \texttt{"Zelcash consensus library"} against an actual build target
named \texttt{libfluxconsensus.\allowbreak{}la} (\shipped, \srccode{fluxd/\allowbreak{}libfluxconsensus.pc.in}{7,9} vs.\
\srccode{fluxd/\allowbreak{}src/\allowbreak{}Makefile.am}{52}); the client name compiled in is the unmodified upstream
\texttt{"MagicBean"} (\shipped, \srccode{fluxd/\allowbreak{}src/\allowbreak{}clientversion.cpp}{17}); and the startup log still
prints \texttt{"Zelcash version \%s (\%s)"} (\shipped, \srccode{fluxd/\allowbreak{}src/\allowbreak{}init.cpp}{918}). These are
stated as discrepancies in an incomplete rebrand, not as deliberate design.

The fork point itself cannot be read off a version string, because \fluxd carries no explicit
record of the upstream \texttt{zcashd} release it branched from. It can nonetheless be bracketed
from consensus evidence rather than guessed. The vendored \texttt{doc/\allowbreak{}zcash-release-notes/}
tops out at the v2.1.0 series (\texttt{release-notes-2.\allowbreak{}1.\allowbreak{}0.\allowbreak{}md} and
\texttt{release-\allowbreak{}notes-\allowbreak{}2.1.0-\allowbreak{}1.\allowbreak{}md}), the former announcing Blossom's mainnet activation
\srcdoc{fluxd/doc/zcash-release-notes/release-notes-2.1.0.md:1-15}; and the network-upgrade
table contains no Heartwood, Canopy, NU5 or Orchard entry
\srccode{fluxd/\allowbreak{}src/\allowbreak{}consensus/\allowbreak{}params.h}{25-38}, \srccode{fluxd/\allowbreak{}src/\allowbreak{}consensus/\allowbreak{}upgrades.cpp}{11-65},
while only the old-style \texttt{librustzcash\_sapling\_*} FFI is called and never
\texttt{orchard} or \texttt{halo2} \srccode{fluxd/\allowbreak{}src/\allowbreak{}main.cpp}{1447}. \Flux therefore derives
from \textbf{\texttt{zcashd} v2.1.0 (Blossom), at or after that release and before v2.1.2
(Heartwood)}. One caution for anyone re-deriving this: the pinned revision in
\texttt{depends/\allowbreak{}packages/\allowbreak{}librustzcash.\allowbreak{}mk} is a 2018 commit, but it is the same commit
\texttt{zcashd} v2.1.0 itself pins, so it corroborates the bracket above rather than dating the fork
earlier \srcbranch{zcash}{v2.1.0}{depends/packages/librustzcash.mk:7}.

Flux's own consensus-upgrade naming diverges completely from Zcash's. \fluxd's \texttt{UpgradeIndex}
enumerates \texttt{BASE\_SPROUT}, \texttt{UPGRADE\_TESTDUMMY}, \texttt{UPGRADE\_LWMA},
\texttt{UPGRADE\_EQUI144\_5}, \texttt{UPGRADE\_ACADIA}, \texttt{UPGRADE\_KAMIOOKA},
\texttt{UPGRADE\_KAMATA}, \texttt{UPGRADE\_FLUX}, \texttt{UPGRADE\_HALVING},
\texttt{UPGRADE\_P2SHNODES}, \texttt{UPGRADE\_PON} (\shipped,
\srccode{fluxd/\allowbreak{}src/\allowbreak{}consensus/\allowbreak{}params.h}{24-39}). There is no dedicated Overwinter- or Sapling-named
entry; Sapling-era rules are instead gated on \texttt{UPGRADE\_ACADIA} (\S\ref{sec:shielded-machinery}
below). \texttt{BASE\_SPROUT} is \texttt{ALWAYS\_ACTIVE} from genesis on mainnet (\shipped,
\srccode{fluxd/\allowbreak{}src/\allowbreak{}chainparams.cpp}{108-110}).

The cryptographic backend is pinned, not vendored fresh. \texttt{librustzcash} is pulled at build time
from \url{https://github.com/zcash/librustzcash/archive/} at git commit
\texttt{06da3b9ac8f278e5\allowbreak{}d4ae13088cf0a4c0\allowbreak{}3d2c13f5}, version \texttt{0.1} (\shipped,
\srccode{fluxd/\allowbreak{}depends/\allowbreak{}packages/\allowbreak{}librustzcash.mk}{2-8}). The Rust crates it depends on are pinned at
their original 2018-era Zcash-project versions: \texttt{bellman} (the Groth16 R1CS prover) at
\texttt{0.1.0} (\shipped, \srccode{fluxd/\allowbreak{}depends/\allowbreak{}packages/\allowbreak{}crate\_bellman.mk}{2-3}) and \texttt{pairing}
(BLS12-381 curve arithmetic) at \texttt{0.14.2} (\shipped,
\srccode{fluxd/\allowbreak{}depends/\allowbreak{}packages/\allowbreak{}crate\_pairing.mk}{2-3}). Nothing in the tree indicates a later
\texttt{librustzcash}, \texttt{bellman}, or \texttt{pairing} upgrade was ever pulled forward.

\subsection{The shielded machinery}
\label{sec:shielded-machinery}

The primitives that make a shielded transaction work — a \note, its commitment $\cm$, its nullifier
$\nf$, the note commitment tree $\tree$ that anchors membership proofs, and the zero-knowledge proof
$\zkproof$ that a spend is valid without revealing which note it spends — are Zerocash-family
constructions~\citep{bensasson2014zerocash} as specified for Zcash's Sprout and Sapling shielded
pools~\citep{zcashprotocol}. \fluxd's implementation lives under \texttt{src/flux/} (namespace
\texttt{libflux}, historically \texttt{libzcash}) plus \texttt{src/\allowbreak{}primitives/\allowbreak{}transaction.\allowbreak{}h}, and it
is compiled into \texttt{libfluxconsensus}/\fluxd for both protocol versions
(\shipped, \srccode{fluxd/\allowbreak{}src/\allowbreak{}Makefile.am}{92-101,542-550}). Table~\ref{tab:shielded-components}
states, per component, whether the code is Zcash's cryptography carried forward unmodified (beyond the
\texttt{libzcash}$\to$\texttt{libflux} rename) or a Flux-added policy layered on top of it; every
Flux-original item in this section is a consensus switch, never a change to the underlying scheme.

\begin{itemize}
\item \textbf{Notes.} \texttt{Note.hpp}\slash\texttt{.cpp} implements \texttt{SproutNote} (fields
  $a_{pk},\rho,r$; $\cm$ computed as $\mathrm{SHA256}(0\mathrm{xb0} \| a_{pk} \| \mathrm{value}_{LE}
  \| \rho \| r)$, \shipped, \srccode{fluxd/\allowbreak{}src/\allowbreak{}flux/\allowbreak{}Note.cpp}{24-41}; $\nf$ via \texttt{PRF\_nf},
  \shipped, \srccode{fluxd/\allowbreak{}src/\allowbreak{}flux/\allowbreak{}Note.cpp}{43-45}) and \texttt{SaplingNote} (fields $d, pk_d, r$;
  $\cm$/$\nf$ delegated to \texttt{librustzcash\_sapling\_compute\_cm}\slash\texttt{\_compute\_nf}, \shipped,
  \srccode{fluxd/\allowbreak{}src/\allowbreak{}flux/\allowbreak{}Note.cpp}{55-93}), with explicit plaintext lead-byte tags \texttt{0x00}
  (Sprout) / \texttt{0x01} (Sapling) (\shipped, \srccode{fluxd/\allowbreak{}src/\allowbreak{}flux/\allowbreak{}Note.hpp}{97-101,156-161}) and
  an outgoing-viewing-key plaintext for sender-side recovery (\shipped,
  \srccode{fluxd/\allowbreak{}src/\allowbreak{}flux/\allowbreak{}Note.hpp}{172-204}). Inherited unmodified.
\item \textbf{Note encryption.} \texttt{Note\allowbreak{}Encryption.\allowbreak{}hpp}\slash\texttt{.cpp} implements
  \texttt{Sapling\allowbreak{}Note\allowbreak{}Encryption} (one-shot, ephemeral $epk$/$esk$; the class comment itself notes it
  is not thread-safe, \srccode{fluxd/\allowbreak{}src/\allowbreak{}flux/\allowbreak{}NoteEncryption.hpp}{30}), trial-decryption routines
  \texttt{Attempt\allowbreak{}Sapling\allowbreak{}Enc\allowbreak{}Decryption}\slash\texttt{Attempt\allowbreak{}Sapling\allowbreak{}Out\allowbreak{}Decryption}, a templated Sprout
  AEAD scheme sized to \texttt{NOTEENCRYPTION\_AUTH\_BYTES}$=16$ bytes, and a
  \texttt{Payment\allowbreak{}Disclosure\allowbreak{}Note\allowbreak{}Decryption} variant used by the payment-disclosure RPCs
  (\S\ref{sec:rpc-surface}) (\shipped, \srccode{fluxd/\allowbreak{}src/\allowbreak{}flux/\allowbreak{}NoteEncryption.hpp}{176-195}).
  Inherited unmodified.
\item \textbf{Note commitment tree ($\tree$).} A generic incremental Merkle tree,
  \texttt{Incremental\allowbreak{}Merkle\allowbreak{}Tree<Depth,Hash>}, instantiated twice: \texttt{Sprout\allowbreak{}Merkle\allowbreak{}Tree} at depth
  29 over SHA-256 compression, and \texttt{Sapling\allowbreak{}Merkle\allowbreak{}Tree} at depth 32 over a Pedersen hash
  (\shipped, \srccode{fluxd/\allowbreak{}src/\allowbreak{}flux/\allowbreak{}IncrementalMerkleTree.hpp}{253-263}). Inherited unmodified.
\item \textbf{Addresses and viewing keys ($\viewkey$).} \texttt{Address.hpp} defines
  \texttt{Sprout\allowbreak{}Payment\allowbreak{}Address}\slash\texttt{ViewingKey}\slash\texttt{SpendingKey} and
  \texttt{Sapling\allowbreak{}Payment\allowbreak{}Address} (diversifier $d$ + $pk_d$), \texttt{Sapling\allowbreak{}Incoming\allowbreak{}Viewing\allowbreak{}Key},
  \texttt{Sapling\allowbreak{}Full\allowbreak{}Viewing\allowbreak{}Key} ($ak,nk,ovk$), and \texttt{Sapling\allowbreak{}Expanded\allowbreak{}Spending\allowbreak{}Key}
  ($ask,nsk,ovk$) (\shipped, \srccode{fluxd/\allowbreak{}src/\allowbreak{}flux/\allowbreak{}Address.hpp}{32-222}).
  \texttt{zip32.h}\slash\texttt{.cpp} implements ZIP-32 hierarchical-deterministic derivation for Sapling
  extended keys (\shipped, \srccode{fluxd/\allowbreak{}src/\allowbreak{}flux/\allowbreak{}zip32.h}{16-18,53-136}). Inherited unmodified.
\item \textbf{Proving system ($\zkproof$).} \texttt{Proof.hpp}\slash\texttt{.cpp} carries the legacy Sprout
  \texttt{PHGRProof} pairing-based proof format (\shipped,
  \srccode{fluxd/\allowbreak{}src/\allowbreak{}flux/\allowbreak{}Proof.hpp}{159-185}) alongside a strict/disabled verifier toggle used at
  reindex time (\shipped, \srccode{fluxd/\allowbreak{}src/\allowbreak{}flux/\allowbreak{}Proof.hpp}{222-229}). \texttt{JoinSplit.hpp} fixes
  \texttt{GROTH\_PROOF\_SIZE}$=48+96+48=192$ bytes — a BLS12-381 compressed $G_1\|G_2\|G_1$ layout —
  and defines \texttt{SproutProof} as a tagged variant of \texttt{PHGRProof} and \texttt{GrothProof}
  (\shipped, \srccode{fluxd/\allowbreak{}src/\allowbreak{}flux/\allowbreak{}JoinSplit.hpp}{22-29}); \texttt{transaction.h} carries the
  Sapling \texttt{Spend\allowbreak{}Description} and \texttt{Output\allowbreak{}Description} wire structures, each holding a
  \texttt{Groth\allowbreak{}Proof zkproof} field (\shipped, \srccode{fluxd/\allowbreak{}src/\allowbreak{}primitives/\allowbreak{}transaction.h}{155-242}).
  Inherited unmodified; see \S\ref{sec:trusted-setup} for the parameters this proving system consumes.
\end{itemize}

\begin{table}[H]
\centering
\small
\begin{tabular}{lll}
\toprule
Component & Provenance & Maturity \\
\midrule
Note ($\note$), commitment ($\cm$), nullifier ($\nf$) & inherited, unmodified (renamed) & \shipped \\
Note encryption & inherited, unmodified (renamed) & \shipped \\
Note commitment tree ($\tree$) & inherited, unmodified (renamed) & \shipped \\
Addresses, viewing keys ($\viewkey$), ZIP-32 & inherited, unmodified (renamed) & \shipped \\
Groth16 proving system ($\zkproof$) & inherited, unmodified (renamed) & \shipped \\
Turnstile enforcement (ZIP-209) & inherited, Flux-disabled & \srccode{fluxd/\allowbreak{}src/\allowbreak{}chainparams.h}{200} \\
Coinbase-must-be-shielded & inherited, Flux-disabled post-height & \srccode{fluxd/\allowbreak{}src/\allowbreak{}main.cpp}{3228-3237} \\
New-inflow restriction at \texttt{UPGRADE\_FLUX} & Flux-original policy & \srccode{fluxd/\allowbreak{}src/\allowbreak{}main.cpp}{1398-1408} \\
Sprout$\to$Sapling migration RPC & inherited, Flux-disabled & \srccode{fluxd/\allowbreak{}src/\allowbreak{}wallet/\allowbreak{}rpcwallet.cpp}{4110-4114} \\
\bottomrule
\end{tabular}
\caption{Provenance of the shielded-value components discussed in this section. ``Flux-disabled''
means the code path is compiled in and reachable but a Flux-added switch keeps it inactive; see
\S\ref{sec:switched-off}.}
\label{tab:shielded-components}
\end{table}

Two shielded protocol versions are active on \fluxd today, and a third is not present at all.
\textbf{Sprout} is active from genesis (\texttt{BASE\_SPROUT} is \texttt{ALWAYS\_ACTIVE}) and accepts
both the legacy \texttt{PHGRProof} and \texttt{GrothProof} encodings, tag-selected per transaction
version (\shipped, \srccode{fluxd/\allowbreak{}src/\allowbreak{}primitives/\allowbreak{}transaction.h}{244-289}). \textbf{Sapling} is gated on
\texttt{Consensus::\allowbreak{}UPGRADE\_ACADIA}, active on mainnet from height 250{,}000 (\shipped,
\srccode{fluxd/\allowbreak{}src/\allowbreak{}chainparams.cpp}{122-125}), with enforcement of the Overwinter flag,
\texttt{SAPLING\_VERSION\_GROUP\_ID}, and the version-4 transaction bounds inside
\texttt{Contextual\allowbreak{}Check\allowbreak{}Transaction} (\shipped, \srccode{fluxd/\allowbreak{}src/\allowbreak{}main.cpp}{1223,1344-1380}). As already
noted, there is no separate Overwinter- or Sapling-named upgrade entry; Flux folded both into
\texttt{ACADIA}. \textbf{Orchard is absent}: no occurrence of the string \texttt{Orchard} appears
anywhere under \texttt{src/} (\shipped, \srccode{fluxd/\allowbreak{}src}{grep -rn Orchard, zero matches}), which is
consistent with the fork point established above — the NU5/Halo2/Orchard-era \texttt{zcashd} code was
never merged. \S\ref{sec:orchard-migration} returns to what adopting it later would and would not
change.

\subsection{The trusted setup, stated exactly}
\label{sec:trusted-setup}

Groth16, as deployed here, requires a one-time per-circuit parameter-generation ceremony whose
output — the proving and verifying keys — must not leak the ceremony's intermediate randomness (the
``toxic waste'')~\citep{groth2016size,bensasson2014zerocash}. \fluxd does not run its own ceremony.

\begin{assumption}[Sapling trusted setup]
\fluxd's Groth16 proving parameters — \texttt{sapling-spend.\allowbreak{}params}, \texttt{sapling-output.\allowbreak{}params},
and \texttt{sprout-groth16.\allowbreak{}params} — are fetched, not generated, by the inherited
\texttt{zcutil/\allowbreak{}fetch-params.\allowbreak{}sh} tool, packaged in the Flux build with a man page still named
\texttt{zcash-\allowbreak{}fetch-\allowbreak{}params} (\shipped, \srccode{fluxd/\allowbreak{}zcutil/\allowbreak{}fetch-params.sh}{5-15}; \srccode{fluxd/\allowbreak{}doc/\allowbreak{}man/\allowbreak{}zcash-fetch-params.1}{4}).
They are re-hosted on a Flux-operated mirror (\texttt{download.\allowbreak{}runonflux.\allowbreak{}io}) and on IPFS, downloaded
in two parts and checked against fixed sha256 pins before use (\shipped,
\srccode{fluxd/\allowbreak{}zcutil/\allowbreak{}fetch-params.sh}{104-141,220-222}). At startup, \texttt{ZC\_Load\allowbreak{}Params()}
requires all three files to exist or calls \texttt{Start\allowbreak{}Shutdown()} (\shipped,
\srccode{fluxd/\allowbreak{}src/\allowbreak{}init.cpp}{754-806,762-775}), then passes independent SHA-512 hash strings for each
file into \texttt{librustzcash\_init\_zksnark\_params()} (\shipped,
\srccode{fluxd/\allowbreak{}src/\allowbreak{}init.cpp}{791-801}). Flux therefore inherits the original \textbf{Zcash Sapling
parameter-generation MPC} — the 2018 ``Powers of Tau''/Sapling ceremony run by the Zcash project
together with outside participants~\citep{bowe2017mpc} — verbatim: it re-hosts the identical
parameter files under their original filenames and role, and performs no additional or independent
ceremony of its own.
\end{assumption}

The claim that the re-hosted files are byte-identical to the ones \texttt{zcashd} itself ships rests
on the filenames, roles, and hash pins matching the well-known public Zcash Sapling parameter set,
together with the crate pins in \S\ref{sec:shielded-lineage} being the original 2018 \texttt{librustzcash}\slash\texttt{bellman}
versions; it was not independently re-derived against a \texttt{zcashd} source tree in the evidence
this section draws on \srcinferred. The paper states it as a strongly supported inference rather than
an independently re-verified fact, and flags the direct hash diff as the bounded check that would
close it.

Table~\ref{tab:trusted-setup} gives the fetch-time integrity pins exactly, since a reviewer checking a
trusted-setup claim will want the exact values rather than a description of them.

\begin{table}[H]
\centering
\small
\begin{tabular}{lll}
\toprule
Parameter file & sha256 (fetch-time pin) & Location \\
\midrule
\texttt{sapling-spend.\allowbreak{}params} & \texttt{8e48ffd2\ldots7efc13} & \srccode{fluxd/\allowbreak{}zcutil/\allowbreak{}fetch-params.sh}{220} \\
\texttt{sapling-output.\allowbreak{}params} & \texttt{2f0ebbcb\ldots f3fb0e4} & \srccode{fluxd/\allowbreak{}zcutil/\allowbreak{}fetch-params.sh}{221} \\
\texttt{sprout-groth16.\allowbreak{}params} & \texttt{b685d700\ldots0ee55b50} & \srccode{fluxd/\allowbreak{}zcutil/\allowbreak{}fetch-params.sh}{222} \\
\bottomrule
\end{tabular}
\caption{Fetch-time sha256 pins for the three Groth16 parameter files \fluxd requires at startup
(hashes elided to their leading/trailing bytes here; full values in
\srccode{fluxd/\allowbreak{}zcutil/\allowbreak{}fetch-params.sh}{220-222}). Runtime loading additionally checks independent
SHA-512 hashes of each file (\srccode{fluxd/\allowbreak{}src/\allowbreak{}init.cpp}{791-801}).}
\label{tab:trusted-setup}
\end{table}

What this assumption buys, and what it costs if it is false, should be stated as plainly as the
assumption itself. If the toxic waste from the original Sapling ceremony was retained by any
participant rather than destroyed, that party could construct false Groth16 proofs that pass
verification while creating shielded value that no spent note ever committed to — that is, forge
notes and mint counterfeit shielded value without breaking any wallet's transparent balance or the
transparent supply at all~\citep{groth2016size,bensasson2014zerocash}. The forgery would be
\emph{undetectable} from inside the shielded pool by construction, because the proof system gives no
signal distinguishing a validly spent note from a forged one; the only external check available is
exactly the pool-balance invariant discussed next, and \S\ref{sec:switched-off} states that this check
is currently disabled on Flux's own chain. The multi-party design of the original ceremony was
specifically intended to make this failure require collusion across many independent
participants~\citep{bowe2017mpc}; Flux's inheritance of the parameters carries that same assumption
forward unchanged, since it is not a party to the ceremony and cannot strengthen or re-verify it after
the fact.

\subsection{What is switched off}
\label{sec:switched-off}

The honest core of this section is that several of the safeguards this machinery was built to
provide are compiled into \fluxd but not active on Flux's live networks.

\begin{description}
\item[The shielded pool has been closed to new entrants since \texttt{UPGRADE\_FLUX}.] This is the
  strongest item in this subsection: the other three restrict or disable a safeguard, this one means
  the pool cannot be funded from the transparent side at all. \texttt{Contextual\allowbreak{}Check\allowbreak{}Transaction}
  rejects any transaction that has transparent inputs (\texttt{tx.vin.size()} nonempty) together with
  a nonempty \texttt{vJoinSplit} or \texttt{v\allowbreak{}Shielded\allowbreak{}Output}, with reject reason
  \texttt{bad-\allowbreak{}txns-\allowbreak{}fluxrebrand-\allowbreak{}vprivacy-\allowbreak{}not-\allowbreak{}empty} (\shipped, \srccode{fluxd/\allowbreak{}src/\allowbreak{}main.cpp}{1398-1408}),
  once \texttt{UPGRADE\_FLUX} is active — mainnet height 835{,}554 (\shipped,
  \srccode{fluxd/\allowbreak{}src/\allowbreak{}chainparams.cpp}{137-139}). The in-source comment states the intent directly:
  ``If flux rebrand is active, we no longer want to allow coins to go into the privacy pool''
  (\srccode{fluxd/\allowbreak{}src/\allowbreak{}main.cpp}{1398-1400}). This is a \textbf{Flux-original consensus policy, not a
  Zcash ZIP} — nothing in the Sprout or Sapling specification restricts shielding transparent value
  this way. The check never inspects \texttt{vShieldedSpend}, so shielded-to-shielded (z$\to$z) and
  shielded-to-transparent (z$\to$t) transactions remain unrestricted; only the transparent-to-shielded
  (t$\to$z) and mixed t/z$\to$z paths that would add new value to the pool are blocked. The pool is
  therefore closed and monotonically non-increasing in value: value can leave for
  the transparent side or move among existing shielded holders, but no new unit can enter.

  At chain tip 2{,}917{,}115 the pool holds 77{,}188.759\,213\,24\flux in Sapling and 19{,}291.177\,350\,41\flux in
  Sprout — 96{,}479.936\,563\,65\flux total (0.0225\% of issued supply) —
  \srcmeasured{\texttt{api.\allowbreak{}runonflux.\allowbreak{}io/\allowbreak{}daemon/\allowbreak{}getblockchaininfo}, field \texttt{valuePools}, at \texttt{blocks}~$=$~2{,}917{,}115}{2026-09-03}.
  \texttt{UPGRADE\_FLUX} activated at height 835{,}554, approximately 10~April~2021 — the calendar
  date is carried only as a source comment, ``Around 10th April 2021''
  (\srccode{fluxd/\allowbreak{}src/\allowbreak{}chainparams.cpp}{138}) — so as of this
  writing (2026-09-03) the pool has been closed to new entrants for approximately 5.4 years. This
  also bounds what the paper can say about shielded privacy in practice: the anonymity set a shielded
  transaction hides within — the notes and holders it is indistinguishable from — cannot grow from new
  participants while the pool stays closed, and neither its note count nor its holder count is exposed
  by any RPC surface this section documents (\S\ref{sec:rpc-surface}), so neither is measured here.
\item[Turnstile enforcement (ZIP-209) is compiled in but dormant.] \texttt{ConnectBlock()} contains
  the inherited check that rejects a block if the running chain-wide Sprout or Sapling balance would
  go negative — the mechanism that would catch exactly the counterfeiting scenario in
  \S\ref{sec:trusted-setup} (\shipped, \srccode{fluxd/\allowbreak{}src/\allowbreak{}main.cpp}{3838-3862}). It is gated by
  \texttt{fZIP209Enabled}, whose class default is \texttt{false} (\shipped,
  \srccode{fluxd/\allowbreak{}src/\allowbreak{}chainparams.h}{200}); the mainnet checkpoint block and Sprout balance-checkpoint
  values that would turn it on are present in the source but left commented out, and the values they
  carry are the unmodified upstream Zcash-mainnet numbers, never adapted to Flux's own chain (\shipped,
  \srccode{fluxd/\allowbreak{}src/\allowbreak{}chainparams.cpp}{322-327}; same pattern on testnet at
  \srccode{fluxd/\allowbreak{}src/\allowbreak{}chainparams.cpp}{527-530}). Only regtest turns it on explicitly (\shipped,
  \srccode{fluxd/\allowbreak{}src/\allowbreak{}chainparams.cpp}{732}). Turnstile enforcement is precisely the mechanism that
  makes the roadmap's own commitment true — ``Supply stays publicly verifiable, with pool totals
  visible so the arithmetic always checks out'' \srcroadmap{flux-roadmap-public.md:174} — and on
  mainnet and testnet it is off. Per G9, the paper states this as the gap between the commitment and
  the deployed system rather than resolving it: turning \texttt{fZIP209Enabled} on, with checkpoint
  values recomputed for Flux's own chain history rather than copied from Zcash's, is what would make
  the sentence true.
\item[Coinbase-must-be-shielded is disabled from \texttt{UPGRADE\_FLUX} onward.] Mainnet and testnet
  set \texttt{f\allowbreak{}Coinbase\allowbreak{}Must\allowbreak{}Be\allowbreak{}Protected = true}; regtest sets it \texttt{false} (\shipped,
  \srccode{fluxd/\allowbreak{}src/\allowbreak{}chainparams.cpp}{89,341,545}). Enforcement inside
  \texttt{Consensus::\allowbreak{}Check\allowbreak{}Tx\allowbreak{}Inputs} is itself gated by whether the ``Flux rebrand'' upgrade is active:
  before height 835{,}554, a coinbase output could only be spent to a shielded address; from that
  height on, the restriction is lifted and coinbase output may be spent transparently (\shipped,
  \srccode{fluxd/\allowbreak{}src/\allowbreak{}main.cpp}{3200,3228-3237}), mirrored on the wallet side (\shipped,
  \srccode{fluxd/\allowbreak{}src/\allowbreak{}wallet/\allowbreak{}wallet.cpp}{3506}). This is a narrower, separate switch from the pool-closure
  restriction discussed first in this subsection: the two share an activation height but not a
  mechanism, and coinbase transparency does not by itself reopen the shielded pool to new value.
\item[\texttt{z\_setmigration} unconditionally throws post-\texttt{UPGRADE\_FLUX}.] The inherited
  ZIP-308 Sprout$\to$Sapling migration RPCs, \texttt{z\_setmigration} and
  \texttt{z\_getmigrationstatus}, are backed by a still-compiled
  \texttt{Async\allowbreak{}RPCOperation\_saplingmigration} (\shipped,
  \srccode{fluxd/\allowbreak{}src/\allowbreak{}wallet/\allowbreak{}asyncrpcoperation\_saplingmigration.h}{11}). But once the Flux-rebrand
  upgrade is active, \texttt{z\_setmigration} immediately throws
  \texttt{"Flux fork active:\allowbreak{} z\_setmigration has been deprecated"} (\shipped,
  \srccode{fluxd/\allowbreak{}src/\allowbreak{}wallet/\allowbreak{}rpcwallet.cpp}{4110-4114}), and \texttt{UPGRADE\_FLUX} activates at mainnet
  height 835{,}554 (\shipped, \srccode{fluxd/\allowbreak{}src/\allowbreak{}chainparams.cpp}{137-139}). The migration operation
  code is reachable only through an internal wallet flag that is itself settable only via this
  now-error-throwing RPC — the migration path has been dead code on mainnet since that height, and
  Sprout value has no in-wallet path to Sapling.
\item[Sapling is gated on \texttt{UPGRADE\_ACADIA}, not a named Overwinter/Sapling upgrade.] As stated
  in \S\ref{sec:shielded-machinery}, this is a naming divergence rather than a functional gap, but it
  means a reader looking for ``Sapling activation'' in \fluxd's upgrade table by Zcash's own naming
  convention will not find it (\shipped, \srccode{fluxd/\allowbreak{}src/\allowbreak{}consensus/\allowbreak{}params.h}{24-39}).
\end{description}

The roadmap's privacy commitments should be read against this closure rather than independently of
it. ``Flux descends from the Zcash lineage. The shielded machinery is part of what we inherited, and
we are bringing private transactions back'' \srcroadmap{flux-roadmap-public.md:167} and ``Shielded by
default in the wallets we make'' \srcroadmap{flux-roadmap-public.md:172} read, taken on their own, like
wallet-integration tasks — pick a construction, ship it in Zelcore and SSP, default new addresses to
\zaddr. Given the closure above, they are not only that. The t$\to$z restriction is a consensus rule
enforced in \texttt{Contextual\allowbreak{}Check\allowbreak{}Transaction}, not a wallet setting, so no wallet change can place
new value in the pool while that rule stands. \textbf{Reopening the pool to new entrants would itself be a
consensus change, and it is not going to happen}: the pools are to be retired instead
(\S\ref{sec:shielded-retirement}, \srcintent{design intake}). An earlier round of
the intake had confirmed reopening as part of a privacy update, and this paper carried it as a
sequenced dependency ahead of wallet integration; that is superseded. The consensus change now in
prospect is the opposite one --- invalidating shielded spends after a sunset height --- and
\S\ref{sec:private-payment} and \S\ref{sec:migration} are updated accordingly. This is a narrower
decision than the choice of successor shielded construction discussed in
\S\ref{sec:orchard-migration}: the roadmap's position that it will not date a cryptographic change
before its audit is scoped concerns which construction to adopt next (Orchard/Halo2 or otherwise), not
the reopening of the already-deployed Sapling and Sprout pools to new transparent value, which needs no
new circuit. It also bears directly on \S21: a private-payment construction can be specified against
the shielded primitives in \S\ref{sec:shielded-machinery}, but with reopening cancelled and both pools
retiring, \S21's constructions are inapplicable to the target architecture, and that section is
retained for its requirement set and its measurements (\Cref{rem:pp-status}).

\begin{figure}[H]
\centering
\begin{tikzpicture}[
  font=\footnotesize,
  pool/.style={draw, rounded corners=3pt, minimum height=15mm, text width=32mm, align=center,
               inner sep=3pt},
  arr/.style={-{Latex[length=2.2mm]}, very thick},
  lbl/.style={font=\tiny, align=center, inner sep=1.5pt},
]
  \node[pool] (t) {\textbf{Transparent pool}\\[1pt]{\scriptsize \taddr}};
  \node[pool, right=44mm of t, fill=black!5] (z)
     {\textbf{Shielded pool}\\[1pt]{\scriptsize Sprout + Sapling, \zaddr}\\[2pt]
      {\tiny \num{96479.93656365}\,\flux}};

  % t -> z : blocked
  %% yshift clears the 8mm blocking bar drawn at the same midpoint, which used to
  %% print straight through this label.
  \draw[arr] ([yshift=4mm]t.east) -- node[lbl, above, pos=0.5, yshift=3.5mm]
      {new inflow \textbf{blocked}\\from height 835{,}554} ([yshift=4mm]z.west);
  \node[draw, very thick, minimum width=3mm, minimum height=8mm, fill=black!25]
      at ($([yshift=4mm]t.east)!0.5!([yshift=4mm]z.west)$) {};

  % z -> t : open
  \draw[arr] ([yshift=-4mm]z.west) -- node[lbl, below, pos=0.5]
      {open} ([yshift=-4mm]t.east);

  % z -> z loop
  \draw[arr] (z.north) to[out=50, in=130, looseness=5]
      node[lbl, above] {\zaddr $\to$ \zaddr: open} (z.north east);

  % turnstile gate, present but inactive
  \node[draw, dashed, rounded corners=2pt, below=13mm of $(t)!0.5!(z)$, text width=48mm,
        align=center, inner sep=3pt] (ts)
      {\scriptsize turnstile check (ZIP-209)\\[1pt]{\tiny compiled in, \texttt{f\allowbreak{}ZIP209Enabled = false}}};
  \draw[dashed, thick] (ts) -- ($(t.east)!0.5!(z.west)$);

  \node[lbl, below=2mm of ts, text width=62mm]
      {The pool can be left and moved within, but not entered. Its value is therefore
       monotonically non-increasing until the rule is lifted.};
\end{tikzpicture}
\caption{Value flow between \fluxd's transparent and shielded pools, showing the turnstile check as
present in code but inactive, and the height-835{,}554 restriction on new transparent-to-shielded
inflow. \shipped, \srccode{fluxd/\allowbreak{}src/\allowbreak{}main.cpp}{1398-1408,3838-3862}; \srccode{fluxd/\allowbreak{}src/\allowbreak{}chainparams.h}{200}.}
\label{fig:shielded-value-flow}
\end{figure}

\subsection{Supply accounting}
\label{sec:supply-accounting}

\fluxd exposes shielded-pool accounting through a \texttt{valuePools} array on the
\texttt{getblockchaininfo} and \texttt{getblockheader} RPCs, with one entry per pool (\texttt{id}
\texttt{"sprout"} or \texttt{"sapling"}), each carrying \texttt{monitored}, \texttt{chainValue},
\texttt{chainValueZat}, \texttt{valueDelta}, and \texttt{valueDeltaZat} (\shipped,
\srccode{fluxd/\allowbreak{}src/\allowbreak{}rpc/\allowbreak{}blockchain.cpp}{84-101}), reported both per block
(\srccode{fluxd/\allowbreak{}src/\allowbreak{}rpc/\allowbreak{}blockchain.cpp}{279-282}) and at the chain tip
(\srccode{fluxd/\allowbreak{}src/\allowbreak{}rpc/\allowbreak{}blockchain.cpp}{1068-1071}). These fields are what a supply audit would read
from: the total issued supply $\supply$ that \S5 defines is, by the roadmap's own description, ``the
transparent set plus the shielded pools'' \srcroadmap{flux-roadmap-public.md:185}, and the roadmap
states that this reconciles with the emission model to within 0.0005\%
\srcroadmap{flux-roadmap-public.md:185}.

The accounting fields are computed and exposed regardless of the turnstile switch discussed in
\S\ref{sec:switched-off}. What the disabled switch removes is not the numbers themselves but the
consensus-enforced invariant that would make a negative pool balance — the signature of exactly the
kind of undetected forgery discussed in \S\ref{sec:trusted-setup} — impossible to sustain on-chain
without triggering block rejection. \texttt{valuePools} is a reporting surface; turnstile enforcement
is the check that would make the report trustworthy on its own. With \texttt{f\allowbreak{}ZIP209Enabled = false}
on both live networks, \fluxd today performs no automated on-chain detection of a negative Sprout or
Sapling pool balance outside of whatever manual monitoring an operator does against the
\texttt{valuePools} field (\shipped, \srccode{fluxd/\allowbreak{}src/\allowbreak{}chainparams.h}{200}). This is the same gap
named in \S\ref{sec:switched-off} and G9, restated here in its supply-accounting form because it is
what connects this section to \S5's supply function: $\supply$ as defined there is only as trustworthy
as the pool-balance fields that compose it, and those fields currently carry no consensus-level
non-negativity guarantee.

\subsection{The RPC surface}
\label{sec:rpc-surface}

The wallet RPC command table exposes the full set of z-address operations
(\shipped, \srccode{fluxd/\allowbreak{}src/\allowbreak{}wallet/\allowbreak{}rpcwallet.cpp}{4994-5021}). This is the toolkit \S21's
private-payment construction is built on, and it is enumerated here precisely rather than by example:

\begin{description}
\item[Key and address management] \texttt{z\_getnewaddress}, \texttt{z\_listaddresses},
  \texttt{z\_exportkey}, \texttt{z\_importkey}, \texttt{z\_exportviewingkey},
  \texttt{z\_importviewingkey} ($\viewkey$), \texttt{z\_exportwallet}, \texttt{z\_importwallet}
  (\shipped, \srccode{fluxd/\allowbreak{}src/\allowbreak{}wallet/\allowbreak{}rpcwallet.cpp}{4994-5018}). Full and incoming Sapling viewing
  keys ($ak,nk,ovk$ and the incoming-viewing-key form) are the objects these last two RPCs move
  (\shipped, \srccode{fluxd/\allowbreak{}src/\allowbreak{}flux/\allowbreak{}Address.hpp}{32-222}), and \texttt{zip32.h} carries the
  ZIP-32 seed fingerprint and parent-FVK tag fields that identify derived keys (\shipped,
  \srccode{fluxd/\allowbreak{}src/\allowbreak{}flux/\allowbreak{}zip32.h}{32,55,101}).
\item[Balance inspection] \texttt{z\_listreceivedbyaddress}, \texttt{z\_listunspent},
  \texttt{z\_getbalance}, \texttt{z\_gettotalbalance} (\shipped,
  \srccode{fluxd/\allowbreak{}src/\allowbreak{}wallet/\allowbreak{}rpcwallet.cpp}{4994-5018}).
\item[Value movement] \texttt{z\_sendmany}, \texttt{z\_shieldcoinbase},
  \texttt{z\_setmigration}\slash\texttt{z\_getmigrationstatus} (dead post-height-835{,}554, \S\ref{sec:switched-off}),
  and \texttt{z\_mergetoaddress} — the last gated behind an explicit \texttt{-\allowbreak{}zmergetoaddress} flag
  (\shipped, \srccode{fluxd/\allowbreak{}src/\allowbreak{}wallet/\allowbreak{}rpcwallet.cpp}{4463-4466}; \srccode{fluxd/\allowbreak{}src/\allowbreak{}init.cpp}{905-906}).
\item[Asynchronous operation tracking] \texttt{z\_getoperationstatus}, \texttt{z\_getoperationresult},
  \texttt{z\_listoperationids} (\shipped, \srccode{fluxd/\allowbreak{}src/\allowbreak{}wallet/\allowbreak{}rpcwallet.cpp}{4994-5018}).
\item[Raw Sprout operations] \texttt{zcbenchmark}, \texttt{zcrawkeygen}, \texttt{zcrawjoinsplit},
  \texttt{zcrawreceive}, \texttt{zcsamplejoinsplit} (\shipped,
  \srccode{fluxd/\allowbreak{}src/\allowbreak{}wallet/\allowbreak{}rpcwallet.cpp}{4994-5021}).
\item[Payment disclosure] \texttt{z\_getpaymentdisclosure} and \texttt{z\_validatepaymentdisclosure}
  are registered under the RPC category \texttt{"disclosure"} (\shipped,
  \srccode{fluxd/\allowbreak{}src/\allowbreak{}wallet/\allowbreak{}rpcwallet.cpp}{5020-5021}), backed by the note-decryption path described
  in \S\ref{sec:shielded-machinery}. They are compiled in but off by default: both require
  \texttt{f\allowbreak{}Experimental\allowbreak{}Mode}, whose global default is \texttt{false} (\inprogress,
  \srccode{fluxd/\allowbreak{}src/\allowbreak{}main.cpp}{77}), set only via \texttt{-\allowbreak{}experimentalfeatures}, plus their own
  per-feature flag; enabling one without the other aborts startup with an explicit error
  (\inprogress, \srccode{fluxd/\allowbreak{}src/\allowbreak{}init.cpp}{895,903-904}). Payment disclosure is present code, not
  an active default.
\end{description}

Payment disclosure and viewing-key export together give a selective-disclosure capability — reveal a
transaction, or a stream of incoming payments, to a chosen third party without revealing anything
else — that is exactly the primitive \S21's payer/network problem needs; \S21 states the construction
it builds from this surface, or states precisely where the surface stops short of it.

\subsection{Post-quantum scope: what is bounded, and why}
\label{sec:pq-scope}

Groth16 as deployed here is a pairing-based SNARK over BLS12-381 — the 192-byte proof format fixed by
\texttt{GROTH\_PROOF\_SIZE} is exactly the compressed $G_1\|G_2\|G_1$ encoding this construction
uses~\citep{groth2016size} (\shipped, \srccode{fluxd/\allowbreak{}src/\allowbreak{}flux/\allowbreak{}JoinSplit.hpp}{22-26}). Its soundness
rests on pairing- and discrete-log-type hardness assumptions over that curve, which a
sufficiently large quantum computer running Shor's algorithm breaks. This is a property of the
scheme itself, not of Flux's use of it, and it holds regardless of who ran the trusted setup in
\S\ref{sec:trusted-setup}: \textbf{the Sapling shielded layer, as deployed, is not post-quantum
secure.} No post-quantum-related code, comment, or TODO exists anywhere under \texttt{src/} —
a repository-wide search for \texttt{quantum} returns zero hits (\shipped, \srccode{fluxd/\allowbreak{}src}{grep -rn quantum, zero matches}).

Given that, the team's stated position is to bound the post-quantum migration deliberately rather
than claim more than the deployed cryptography supports. What is to become post-quantum is
consensus, transparent transactions, and the keys securing them; the shielded layer is explicitly
placed out of scope for now, with post-quantum privacy construction to follow later as the field
matures \srcintent{design intake} (\vision; after the retirement decision of
\S\ref{sec:shielded-retirement}, any such construction is a future option, not a scheduled item).
This scope decision is presented in the design intake
as coherent and defensible rather than as an admitted gap, on the reasoning that privacy is already
optional at the protocol level (the roadmap's own position, quoted at the head of this section), so a
transparent-layer-first migration does not leave the
part of the system that everyone must use exposed while the part that is opt-in waits. The paper
adopts that framing because the alternative is worse: an unqualified ``quantum-ready'' claim made
while a live Groth16-based shielded pool sits underneath it would be a claim a reviewer familiar with
the scheme would immediately falsify. A scope statement this specific — consensus and transparent
keys first, shielded privacy deferred and named as deferred — is the position that survives that
scrutiny; a blanket claim is not.

One distinction has to be kept strict, because it is the likeliest error in this area. Removing a
trusted setup and achieving post-quantum security are two different properties, and a future
construction can have either without the other. Halo2, the proof system behind Zcash's Orchard
pool, eliminates the per-circuit trusted setup entirely by replacing it with a transparent, universal
setup~\citep{bowe2019halo} — but it does so by resting on discrete-log hardness over a different
curve, which is exactly as vulnerable to a quantum adversary as the pairing assumption Groth16 uses.
Adopting Halo2/Orchard, discussed next, would answer the trusted-setup question and would do nothing
for the post-quantum question. The two must not be conflated in either direction.

\subsection{What a migration to Orchard/Halo2 would, and would not, buy}
\label{sec:orchard-migration}

Orchard is absent from \fluxd today (\S\ref{sec:shielded-machinery}), and no committed timeline or
specific target construction exists yet: the roadmap places ``Chain privacy implementation'' in H2
2027, explicitly ``on the schedule set by the Q4 2026 plan''
\srcroadmap{flux-roadmap-public.md:158} (\planned), and states that the Q4 2026 privacy groundwork
work is where ``which shielded construction to adopt'' is still to be settled
\srcroadmap{flux-roadmap-public.md:177} (\planned). Whether that construction turns out to be
Orchard/Halo2 specifically, some other Zcash-lineage successor, or something else is not yet decided
by the team's own account, so this subsection is stated as \vision rather than as a scheduled item.

The argument for making that move is specific, not generic, and worth stating precisely: adopting a
Halo2-class proving system removes the inherited trusted setup described in
\S\ref{sec:trusted-setup} entirely, and that directly serves the stated decentralization requirement
— a chain whose privacy layer's soundness depends on a multi-party ceremony run once, by a fixed set
of participants, in 2018, is not decentralized in the same sense the rest of the design intake asks
for \srcintent{design intake} (\vision). That argument stands on its own regardless of the
post-quantum question, and the paper states it as such rather than bundling it with a PQ claim it
does not support (\S\ref{sec:pq-scope}).

What it would not buy is equally specific: no candidate successor construction under discussion is
post-quantum secure, so a migration to Orchard/Halo2 closes the trusted-setup gap while leaving the
shielded layer's long-term cryptographic exposure to a quantum adversary exactly where it is today.
The team's stated position on timing is to not pin one: ``We will not pin a date on a cryptographic
change before its audit is scoped. The direction is not in question''
\srcroadmap{flux-roadmap-public.md:177} (\planned). The paper records that position rather than
supplying a date the source does not give.

\subsection{Retiring the pools: the decision, and the measurements behind it}
\label{sec:shielded-retirement}

\textbf{Both shielded pools are to be retired} \srcintent{design intake}
(\planned). The team raised the possibility alongside reopening and migrating --- destroying the
existing pool outright to simplify the chain, reduce maintained code and obtain a more thorough audit
of what remains, with a stated preference for C++ over Rust
\srcintent{design intake} --- and, presented with the analysis below, took it.

This subsection records the measurement that motivated the decision, the constraint that made it
sharper than it first appears, and the schedule. The argument is retained in full rather than
compressed to its outcome, because a decision to destroy a live mechanism holding user value should
remain auditable against the reasoning that produced it.

\paragraph{What is actually in the pools.} Both legacy pools are live and both are small.

\begin{table}[H]
\centering\footnotesize
\caption{Shielded pool balances at height \num{2917115}, from \fluxd's own \texttt{valuePools}
reporting surface. Circulating supply is the emission model's value at that height
(\Cref{sec:monetary-policy}). \shipped
\srcmeasured{api.runonflux.io/daemon/getblockchaininfo, field \texttt{valuePools}}{2026-09-03}.}
\label{tab:pool-balances}
\begin{tabular}{@{}lrr@{}}
\toprule
Pool & Balance (\flux) & Share of circulating supply \\
\midrule
Sprout (JoinSplit lineage) & \num{19291.17735041} & \SI{0.0045}{\percent} \\
Sapling & \num{77188.75921324} & \SI{0.0180}{\percent} \\
\midrule
Total shielded & \textbf{\num{96479.93656365}} & \textbf{\SI{0.0225}{\percent}} \\
\bottomrule
\end{tabular}
\end{table}

\noindent Three facts follow, and together they carry the argument. First, the combined shielded
balance is about one part in \num{4500} of circulating supply. Second, the balance is
\emph{monotonically non-increasing}: the transparent-input gate of
\S\ref{sec:switched-off} rejects any transaction carrying both a transparent input and a
shielded output, so value has been unable to enter either pool since height \num{835554}
\srccode{fluxd/src/main.cpp}{1398--1408}. Third --- and this is the fact most often missed --- the
gate keys on \texttt{tx.vin.size()} alone, so \zaddr${}\to{}$\zaddr and
\zaddr${}\to{}$\taddr transactions remain fully valid. The pools are sealed from outside and
unimpaired within: existing holders retain full use and can exit at will.

\paragraph{The constraint that decides it.} The stated preferences --- newest Zcash work, less code,
a more thorough audit, and C++ over Rust --- cannot all be satisfied, because they point in opposite
directions. Sapling proving and verification in \fluxd run through \texttt{librustzcash}
(\S\ref{sec:shielded-machinery}), which is Rust. Orchard's Halo2 proving system exists as a Rust
implementation and has no production C++ equivalent. Zcash's own trajectory has been to move
consensus-critical cryptography \emph{into} Rust, not out of it. So ``port the newest Zcash work''
and ``prefer C++'' are not jointly achievable: adopting newer Zcash cryptography means more Rust in
the consensus path, not less \srcinferred.

That leaves exactly three self-consistent positions, and the preference set selects among them
cleanly:

\begin{table}[H]
\centering\footnotesize
\caption{The three coherent options, against the team's own stated preferences.}
\label{tab:shielded-options}
\begin{tabular}{@{}lllll@{}}
\toprule
Option & Chain privacy & Rust in consensus & Trusted setup & Audit surface \\
\midrule
(A) Reopen Sapling & yes & retained & inherited, two ceremonies & unchanged, large \\
(B) Retire both pools & \textbf{none} & \textbf{removable} & \textbf{none} & \textbf{smallest} \\
(C) Retire, later adopt Orchard & deferred & increased & none (Halo2) & large, later \\
\bottomrule
\end{tabular}
\end{table}

\paragraph{Option (B) is the decision; (C) remains available later.} Four arguments, in order of
weight.

\begin{enumerate}
  \item \textbf{A shielded pool does not buy what the thesis needs.} The workload is public by
  construction. A Flux application specification is a permanent on-chain message, and the
  control plane serves the global specification set over a public endpoint
  \srccode{flux-domain-manager/src/services/flux/dataFetcher.js}{47-71}. What a tenant deploys, the
  resources it claims and the owner it binds to are therefore public whatever the payment rail does.
  Shielding the payment hides the \emph{funding source}, not the workload graph. That is a real but
  narrow benefit, and \S\ref{sec:private-payment} shows it is also the harder half to construct
  correctly.
  \item \textbf{Every stated preference points the same way.} (A) and (C) both keep or expand Rust in
  the consensus path; only (B) permits removing \texttt{librustzcash}, the JoinSplit and Sapling
  circuits, the note-commitment trees and the nullifier sets from the daemon. If a smaller,
  C++-centred, thoroughly audited consensus layer is the goal, (B) is the only option that delivers
  it.
  \item \textbf{It closes a supply-accounting gap the paper already documents.}
  \Cref{sec:supply-accounting} records that \texttt{valuePools} is a reporting surface with no
  consensus enforcement behind it: a negative pool balance --- the signature of an inflation bug in
  the shielded code --- would not by itself cause block rejection, because ZIP-209 turnstile
  enforcement is compiled in but dormant against unmodified upstream-Zcash checkpoint constants
  \srccode{fluxd/src/chainparams.cpp}{322-327}. Removing the pools removes that entire class of
  unauditable supply risk, and makes total supply exactly reconstructible from transparent UTXOs.
  For a design whose monetary section rests on supply transparency, that is a direct gain.
  \item \textbf{The value at risk is bounded and measured.} \num{96479.94}\flux, in a pool that
  cannot grow. This is the number the decision should be made against, and it is small enough that
  a generous migration window costs little and dishonours no one.
\end{enumerate}

\paragraph{The retirement schedule.} The decision is settled; the heights below were the paper's
recommendation and are confirmed by the author (Remark~\ref{rem:shield-schedule-confirmed}).

\begin{itemize}
  \item \textbf{Announcement at publication}, i.e. with this paper (2026-09-07), alongside the notice
  of the parallel-asset one-way switch at $\hopen$ (\S\ref{sec:pa-sequencing}). One consolidation
  message covering both is better than two, and the holder sets barely overlap.
  \item \textbf{Sunset at $\hshield = \num{3600000}$}, about 2027-04-27 at the \SI{30}{\second}
  target: \num{232} days after publication, \num{150000} blocks after $\hopen$ and \num{187502}
  blocks \emph{before} $\hpa$ \srcintent{08-answers-round-2.md, round 37}. After this height,
  transactions carrying \texttt{vShieldedSpend} or \texttt{vJoinSplit} are invalid; notes not swept
  are permanently unspendable. Seven and a half months is deliberately longer than the roughly
  five-month parallel-asset window (\S\ref{sec:bridging-migration}): that window ends in forfeiture
  to a named recipient, whereas this one ends in destruction, which warrants more notice for less
  value. The height is a consensus rule and must ship in a \texttt{fluxd} release ahead of it.
  \proposednote{Carry it in the \texttt{fluxd} upgrade the roadmap already plans for H1 2027
  (collateral maturity, \S\ref{sec:migration}), so that retirement adds no release of its own before
  \PNR{}.}
  \item \textbf{No change to exit paths during the window.} \zaddr${}\to{}$\taddr already works and
  needs no consensus change to keep working; the window requires an announcement, not a code change.
  \item \textbf{Activate ZIP-209 turnstile enforcement with Flux-correct constants for the duration
  of the window.} Whatever is decided about the pools' future, the current state --- enforcement
  code compiled in but dormant, parameterised by another chain's checkpoints --- is the worst of the
  available positions, because it carries the maintenance cost of the mechanism without its
  guarantee. During a drain-only window the turnstile is a cheap and exactly-correct safety net: the
  pool balance must only decrease, which is precisely what the turnstile checks.
  \item \textbf{Direct outreach, not announcement alone.} The holder set is small enough to make
  this tractable and small enough that a purely passive notice would be a poor effort.
\end{itemize}

\begin{remark}[Schedule confirmed]
\label{rem:shield-schedule-confirmed}
The heights above were proposed by this paper and are \textbf{confirmed by the author}
\srcintent{design intake}; the sunset was then brought forward to $\hshield = \num{3600000}$ and the
announcement to publication \srcintent{08-answers-round-2.md, round 37}. Activating ZIP-209 turnstile enforcement with Flux-correct constants for
the window remains this paper's recommendation, not a confirmed decision: round 11 confirmed the
numbers, and a switch is not a number \srcintent{design intake}. The retirement
decision and its heights are settled, and \S\ref{sec:private-payment}'s construction questions are
consequently not live (\Cref{rem:pp-conditional}).
\end{remark}

\begin{remark}[What retirement costs, stated without softening]
\label{rem:retirement-cost}
Three things are given up, and the paper names them rather than letting the argument above stand
unopposed.

\textbf{Flux will have no chain-level privacy.} Not reduced privacy --- none. Every transfer, every
balance and every payment for a deployment will be publicly linkable for anyone willing to do the
analysis, exactly as on Bitcoin. The design intake asked for privacy-preserving settlement rails
\srcintent{design intake}; after retirement the honest description is that Flux settles
transparently and treats payment privacy as an unsolved problem rather than a delivered feature.
\Cref{sec:pq-scope} and \S\ref{sec:private-payment} are the two places this paper previously
claimed otherwise, and both are now corrected.

\textbf{Value will be destroyed at $\hshield$.} Notes not swept before the sunset become
permanently unspendable. The bound is \num{96479.94}\flux{} and shrinking, and the window is long,
but this is confiscation by inaction and calling it anything softer would be dishonest. It is the
reason the schedule recommends direct outreach rather than announcement alone.

\textbf{The option is one-directional in practice.} Rebuilding a shielded pool later means adopting
Orchard/Halo2 or an equivalent --- a multi-year cryptographic project, in Rust, against the stated
language preference. Option~(C) remains formally available and should not be described as though it
were cheap.
\end{remark}

%% Drain this section's float queue before the next begins.
\FloatBarrier
